\documentclass[11pt]{article}
\usepackage[margin=1in]{geometry}
\usepackage{amsmath, amssymb}
\usepackage{enumitem}

\title{Average Log-Likelihood Ratio: Formulas}
\author{}
\date{}

\begin{document}
\maketitle

\section{Setup}

A position frequency matrix (PFM) has 4 rows (A, C, G, T) and $W$ columns
(positions).  Column $j$ is a probability vector
$\mathbf{f}_j = (f_{A,j},\, f_{C,j},\, f_{G,j},\, f_{T,j})$
with $\sum_i f_{i,j} = 1$.

Let $\mathbf{b} = (b_A, b_C, b_G, b_T)$ be a background distribution
(default: uniform, $b_i = 0.25$).

\section{Column-level ALLR}

Given two columns $C_1$ and $C_2$ with frequency vectors
$\mathbf{f}_1$ and $\mathbf{f}_2$:
\begin{equation}
  \mathrm{ALLR}(C_1, C_2)
  = \frac{1}{2}\!\left[
      \sum_{i \in \{A,C,G,T\}} f_{i,1} \ln\!\frac{f_{i,2}}{b_i}
    + \sum_{i \in \{A,C,G,T\}} f_{i,2} \ln\!\frac{f_{i,1}}{b_i}
  \right].
\end{equation}
The first sum asks how well $C_2$ explains observations drawn from $C_1$
(relative to background); the second sum is the reverse.
Averaging makes the measure symmetric.

\medskip\noindent
Properties:
\begin{itemize}[nosep]
  \item $\mathrm{ALLR}(C_1, C_2) = \mathrm{ALLR}(C_2, C_1)$.
  \item If both columns equal the background, $\mathrm{ALLR} = 0$.
  \item Identical informative columns yield a positive score;
        dissimilar columns yield a low or negative score.
\end{itemize}

\section{Matrix-level ALLR with sliding alignment}

Let PFM $A$ have width $w$ and PFM $B$ have width $W \ge w$.
For offset $k = 0, 1, \dots, W - w$, define the aligned score:
\begin{equation}
  S(k) = \frac{1}{w} \sum_{j=1}^{w}
         \mathrm{ALLR}\!\bigl(A_j,\; B_{k+j}\bigr),
\end{equation}
where $A_j$ is column $j$ of $A$ and $B_{k+j}$ is column $k+j$ of $B$.

The best-match score and offset are:
\begin{equation}
  S^* = \max_{k} S(k), \qquad k^* = \arg\max_{k} S(k).
\end{equation}

If $A$ is wider than $B$, their roles are swapped so the shorter matrix
always slides over the longer one.

\section{P-value by permutation test}

The null hypothesis is: the observed similarity $S^*$ is due to chance.

\subsection{Procedure}

\begin{enumerate}[nosep]
  \item Generate $N$ random PFMs of the same width as the target.
        Each column is drawn independently from
        $\mathrm{Dirichlet}(\alpha_0\, \mathbf{b})$,
        where $\alpha_0 > 0$ is a concentration parameter
        (default: $\alpha_0 = 4$).
  \item For each random PFM $R_n$, compute $S_n^*$ by sliding alignment
        against the input PFM (same procedure as above).
  \item The p-value is:
\end{enumerate}
\begin{equation}
  P = \frac{\#\{n : S_n^* \ge S^*\} + 1}{N + 1}.
\end{equation}
The $+1$ in numerator and denominator is the standard pseudo-count
correction that avoids $P = 0$ and accounts for the observed score
itself being one draw from the distribution.

\subsection{Interpretation}

Small $P$ (e.g.\ $P < 0.05$) indicates the match is unlikely under the
null model.  The permutation approach automatically accounts for motif
width and information content without parametric assumptions.

\end{document}
